% Sample document for the README editor preview; transcribed from the bundled
% set-theory photo and given a title for the scripted editing demonstration.
\documentclass[11pt]{article}
\usepackage[margin=1in]{geometry}
\usepackage{amsmath,amssymb}
\title{Sets --- Chapter 1}
\author{}
\date{}
\begin{document}
\maketitle

\noindent A \textbf{set} is any well-defined collection of elements or members.

\section*{Describing a set}
List the elements between braces (a finite number of elements), or specify a
property that the elements of the set have in common:
\begin{align*}
  R &= \{1,2,3,4\} \\
  R &= \{x : x \text{ is a real number}\}.
\end{align*}
The second description is the \textbf{rule method}, used for infinite sets.

\section*{Order and repetition}
The order of elements in a set is not important:
\[
  \{1,2,3\} = \{2,1,3\} = \{2,3,1\}.
\]
Repeated elements are ignored:
\[
  \{1,2,3,1\} = \{1,2,3\}.
\]

\section*{Notation}
Uppercase letters denote sets: $A,B,C,X,\ldots$.
Lowercase letters denote elements: $a,b,c,x,\ldots$.
\begin{align*}
  x \in A &\qquad x \text{ belongs to the set } A, \\
  x \notin A &\qquad x \text{ does not belong to the set } A.
\end{align*}

\section*{Fuzzy set}
A set that is not well-defined (random elements put together), for example:
\[
  X = \left\{0,4,-2,7.5,-\frac{22}{7}\right\}.
\]

\section*{Venn diagrams}
Sets are represented as geometrical shapes for easier representation. Example:
\[
  A = \{1,2,3,4\}, \qquad B = \{3,4,5,6,7\}.
\]
\end{document}
